\documentclass[10pt,letterpaper,final]{report}
\usepackage[margin=1in]{geometry}
\usepackage[utf8]{inputenc}
\usepackage{amsmath}
\usepackage{amsfonts}
\usepackage{amssymb}
\usepackage{amsthm}
\renewcommand\qedsymbol{$\blacksquare$}
\usepackage{enumerate}
\usepackage{hyperref}
\usepackage{pdfpages}
\usepackage{graphics}
\usepackage{graphicx}
\usepackage{tikz}
\usepackage{tikz-qtree}
\usetikzlibrary{automata,arrows}
\usepackage{mdframed}
\usepackage[
  backend=biber,
  style=numeric-comp,
  sorting=none
]{biblatex}
\addbibresource{references.bib}

\usepackage{minted}
\usemintedstyle{colorful}
\usepackage{xltabular}
\usepackage{pdflscape}

% Based on template from COSC-417 at Towson University, originally authored by Marius Zimand

\renewcommand{\thesection}{\arabic{section}}

\begin{document}
\begin{center}
\fbox{
  \vbox{
    \begin{flushleft}
      Jonathan Llovet \\  % authors' names
      EN.605.202.81 \\  %class
      2026-06-23 \\  % date
      Lab 1 \\ % ID
    \end{flushleft}
    \center{\Large{\textbf{Use of Stacks - Converting Prefix to Postfix Expressions}}}
    %\end{mdframed}
  } % end vbox
} % end fbox
\vline
\end{center}

\noindent\textbf{Prompt:}

In this assignment you will a stack to convert prefix expressions directly to postfix expressions. You may not use recursion as we will use recursion to revisit this problem in Lab 2.


Write a program that accepts a prefix expression containing single letter operands and the operators \verb|+, -, *, /, $| (\verb|$|representing exponentiation). Output the corresponding postfix epression.

\begin{itemize}
  \item For example, if your input is \verb|*AB| then output should be \verb|AB*|. Output of \verb|BA*| is considered incorrect.
\end{itemize}

In your analysis, be sure to discuss the implementation you choose and why, why a stack makes sense. Consider a recursive solution (you should not implement recursion) and compare it to your iterative solution. Is one better than the other? Why? Tell us what you learned and what you would do differently. Be sure to review the Programming Assignment Guidelines for specific requirements for the Analysis and before submitting this assignment.

Don't forget to do reasonable error checking and try to make your error messages specific and helpful.  You may not use library functions. You must write your own code and, specifically, you must write the stack code. Be sure to include the stack source code in your submission. You must read and write from named files.

The required input is provided in a separate text file.

In processing the input, you should expect to read line by line and within a line, character by character, as specified in the programming assignment guidelines. This will allow you to parse the expression as you read it.

\tableofcontents

\pagebreak

\section{Analysis}

This lab illustrates the use of stacks in the conversion of prefix expressions to postfix expressions. The prefix expressions will be read in from a file and then output to a file, both of which are provided at the command line. Errors will be printed to stderr for review. The lab is meant to be run as a module. See \texttt{README.md} for running instructions.
We want to convert expressions in this form: \texttt{+AB} to this form \texttt{AB+}. Both of these are equivalent to the probably more familiar infix notation for the same expression \texttt{A+B}. Put another way, we want to convert from Polish notation to Reverse Polish notation. \parencite{EnWikipediaorgReversePolishNotation}

TODO: Add discussion of recursive algorithm.
TODO: Add discussion of stack and ADTs.
TODO: Add discussion of algorithmic complexity
TODO: Add discussion of design choices


\section{Strategy}

TODO: Talk about pipeline, concerns about equivalence, in a subsequent version I should write a reverse converter or just refactor the existing one to handle both directions, since that would sufficient for isomorphism

TODO: Discuss interpreter as heuristic to help ensure correctness, especially of more complex expressions

TODO: Make diagram of the core flow of the program

\section{Structure}

\begin{verbatim}
lab1
|-- __init__.py            | Python-required for module
|-- __main__.py            | Main entrypoint
|-- convert                | Package that handles actual pre->post conversion
|   |-- __init__.py        | Python-required for package
|   `-- converter.py       | The pre->post conversion pipeline
|-- lab1.py                | IO entrypoint for processing files and printing errors
|-- launch.json            | Configuration for IDE
|-- parse                  | Package to handle parsing expressions
|   |-- __init__.py        | Python-required for package
|   |-- evaluate.py        | Evaluates expressions to help with validity checks
|   |-- parser.py          | Parsing for expressions, symbol lookup, and 
|   `-- validate.py        | Validation functions that check expressions, symbols
|-- stack                  | Package for an array-based stack
|   |-- __init__.py        | Python-required for package
|   `-- stack.py           | Implementation of an array-based stack (uses Python list)
`-- tests                  | "Testing, 1, 2"
    |-- __init__.py        | Python-required for package
    |-- test_converter.py  | Tests for the convert package
    |-- test_parser.py     | Tests for the parse package
    |-- test_postfix.py    | Tests for postfix operations
    |-- test_prefix.py     | Tests for prefix operations
    |-- test_stack.py      | Tests for the stack package
    `-- test_validate.py   | Tests for the validate package

5 directories, 19 files
\end{verbatim}

\section{Algorithms}

TODO: Add discussion of stack algorithms from class slides and their generic forms

\section{Testing}

Note: I used test-driven development with unit tests written in \texttt{unittest} to help with the development of my code and to ensure that I did not break portions of the codebase as I built new features, performed refactors, and did other chores. I referenced several resources to get back up to speed with testing in python using \texttt{unittest}. \parencite{PythonDocumentationUnittestUnit, CodewithJosh2024HowToTest, CoreySchafer2017PythonTutorial, ArjanCodes2025HowToWriteGreatUnitTestsInPython}
I checked multiple resources for information about error handling \parencite{PythonDocumentationErrorsAndExceptions}, \texttt{match} syntax \parencite{PythonEnhancementProposalsPEP0636}, and typing \parencite{PythonDocumentationTypingLiterals}. In addition to these, I referenced the Google Python Style Guide \parencite{GooglePythonStyleGuide} for conventions about docstring formatting.

\section{Demo}

For the following, I assume that I start in the root of the Github repository. From some appropriate location on your filesystem, clone the repository and move into the root of the repo.
{\small
\begin{verbatim}
  git clone https://github.com/jllovet/jhu-en-605-202-su26-data-structures.git
  cd jhu-en-605-202-su26-data-structures
\end{verbatim}
}

\noindent From there, move into the folder for this lab.

{\small
\begin{verbatim}
  cd lab-1
\end{verbatim}
}

\noindent From here, the python module can be run as follows.

{\small
\begin{verbatim}
  python -m lab1 resources/input/in.txt resources/output/output.txt
\end{verbatim}
}

\noindent The input and resulting output are shown below:

\begin{center}
\begin{minipage}[t]{0.46\textwidth}
  \centering
  \textbf{\texttt{resources/input/in.txt}}\\
  \inputminted[frame=single,fontsize=\footnotesize]{text}{resources/input/in.txt}
\end{minipage}\hfill
\begin{minipage}[t]{0.46\textwidth}
  \centering
  \textbf{\texttt{resources/output/output.txt}}\\
  \inputminted[frame=single,fontsize=\footnotesize]{text}{resources/output/output.txt}
\end{minipage}
\end{center}

\noindent When the program is executed, the error message on the following page will be written to \texttt{stderr}:

\pagebreak

{\small
\begin{center}
\begin{verbatim}
WARNING: 19 errors found during conversion
WARNING: Invalid expressions not written to output file

Check expressions to ensure the are well formed and only use allowed symbols.

Note: whitespace is not allowed. Expressions with whitespace will be rejected.

Allowed symbols: ABCDEFGHIJKLMNOPQRSTUVWXYZ+-*/$()

Example valid expressions:
	(base case):	'A'	-> 'A'
	(single op):	'+AB'	-> 'AB+'
	(more ops):	'-+ABC'	-> 'AB+C-'
--------------------------------------------------------------------------------
ERRORS:
resources/input/in.txt - line 2: 'AB' is not a valid prefix expression
resources/input/in.txt - line 3: 'AB+' is not a valid prefix expression
resources/input/in.txt - line 6: '++AB' is not a valid prefix expression
resources/input/in.txt - line 7: '+++' is not a valid prefix expression
resources/input/in.txt - line 8: 'A+B' is not a valid prefix expression
resources/input/in.txt - line 9: 'A-B' is not a valid prefix expression
resources/input/in.txt - line 10: 'A=B' is not a valid prefix expression
resources/input/in.txt - line 11: 'A#B' is not a valid prefix expression
resources/input/in.txt - line 12: 'A@B' is not a valid prefix expression
resources/input/in.txt - line 13: 'A$B' is not a valid prefix expression
resources/input/in.txt - line 17: 'miscellaneous' is not a valid prefix expression
resources/input/in.txt - line 18: '+La-sc*ia$te' is not a valid prefix expression
resources/input/in.txt - line 19: '$+' is not a valid prefix expression
resources/input/in.txt - line 23: '/A+BC +C*BA  ' is not a valid prefix expression
resources/input/in.txt - line 24: '*-*-ABC+BA  ' is not a valid prefix expression
resources/input/in.txt - line 25: '/+/A-BC-BA  ' is not a valid prefix expression
resources/input/in.txt - line 26: '*$A+BC+C-BA ' is not a valid prefix expression
resources/input/in.txt - line 27: '//A+B0-C+BA' is not a valid prefix expression
resources/input/in.txt - line 28: '*$A^BC+C-BA' is not a valid prefix expression
\end{verbatim}
\end{center}
}

\pagebreak

\section{Source}


\pagebreak

\section{Source Code}
The full source code, including LaTeX, tooling, other artifacts, and git history is available on my Github repository for the class here:
\begin{center}
  \url{https://github.com/jllovet/jhu-en-605-202-su26-data-structures}
\end{center}

\printbibliography[heading=bibintoc]

\end{document}